% ======================================================================
%  CAPÍTULO 17 — PREGUNTAS FRECUENTES
% ======================================================================

\part{Cierre}\label{part:cierre}

\chapter{Preguntas frecuentes}\label{ch:faq}

\begin{roadmap}
Dieciséis capítulos de teoría, tres laboratorios y dos años de
literatura convulsa dejan, al final, las preguntas de siempre: ¿qué
es esto?, ¿por qué así?, ¿está roto?, ¿lo uso? Este capítulo las
responde en formato breve y con punteros exactos a las partes donde
cada respuesta está desarrollada. Si solo se va a leer una sección
del libro después de la introducción, que sea esta; si algo de aquí
resulta sorprendente o insuficiente, los punteros llevan al capítulo
correspondiente.
\end{roadmap}

\section{El esquema, en preguntas cortas}

\begin{enumerate}[leftmargin=1.9em,itemsep=5pt]
  \item \textbf{¿Qué es Classic McEliece, en una frase?} Un KEM
        post-cuántico: el criptosistema de McEliece (1978) en la
        forma dual de Niederreiter, con códigos de Goppa binarios,
        estandarizado por ISO en 2026. Es famoso por dos extremos:
        las claves públicas más grandes de la criptografía
        moderna---de $261\,120$ a $1\,357\,824$ bytes---y algunos de
        los criptogramas más pequeños ($96$--$208$ bytes). Su
        seguridad se ancla en la dureza conjetural de la
        decodificación por síndromes, un problema bajo escrutinio
        desde 1962~\cite{cmspec,cmsecurity}.

  \item \textbf{¿Es «el McEliece de 1978» o es otra cosa?} El mismo
        esqueleto (códigos de Goppa, trampa de Patterson), con tres
        piezas reingenierizadas: la forma dual, donde el mensaje es
        un vector de error de peso fijo y no un texto; la forma
        sistemática con semilla de 32 bytes, que reduce el
        almacenamiento y hace reproducible la generación
        (capítulo~\ref{ch:keygen}); y el transformado a KEM con
        rechazo implícito (capítulos~\ref{ch:kem}
        y~\ref{ch:encapdecap}). La trampa---decodificar es fácil con
        $g$ y $\alpha$ y es el problema duro sin
        ellos---es exactamente la de~\cite{mceliece1978}.

  \item \textbf{¿Por qué las claves públicas son gigantescas?}
        Porque la clave pública es, precisamente, una matriz binaria
        aleatoria de $mt\times k$ bits (la $T$ de
        $(\mathrm{I}_{mt}|T)$). No hay anillo ni estructura que la
        comprima, y esa ausencia de estructura es parte del argumento
        de seguridad: menos álgebra visible, menos superficie para el
        criptoanálisis estructural (capítulo~\ref{ch:estructura}).
        El contraste está en la tabla~\ref{tab:param-contexto}: unas
        $200$ veces las claves de ML-KEM; a cambio, criptogramas
        entre $5$ y $8$ veces más pequeños~\cite{cmsecurity}.

  \item \textbf{¿Cifra mensajes o claves?} Es un KEM: acuerda una
        clave de sesión de $32$ bytes, que después cifra el mensaje
        en un esquema simétrico (AES-GCM, ChaCha20). En un protocolo
        real se combina en híbrido con un intercambio clásico
        (X25519): si un lado cae, el otro sostiene; los despliegues
        conocidos usan todos el modo
        híbrido~\cite{cmimpl,bernstein2026}.

  \item \textbf{¿Qué tan rápido es?} En los términos de conteo del
        capítulo~\ref{ch:encapdecap}: la encapsulación es una
        multiplicación matriz-vector sobre la clave pública
        (XORs; $\sim2^{23{,}4}$ operaciones de bit en el conjunto
        mayor) más los hashes; la desencapsulación añade Patterson.
        Lo caro es la \emph{generación} (reintentos de
        $\textsf{MatGen}$: $\sim3{,}4$ de media en la forma base,
        una sola en las variantes \texttt{f}), pero ocurre una vez
        por clave~\cite{cmrationale}.

  \item \textbf{¿Por qué el «mensaje» es un vector de error?} Porque
        el problema duro de un código es el síndrome: dada
        $C=He$, recuperar $e$ es exactamente
        decodificar (capítulo~\ref{ch:isd}). La especificación no
        disimula el corazón del problema: se lo pide al atacante en
        su forma desnuda, y restaura la uniformidad con el
        transformado a KEM (capítulos~\ref{ch:kem}
        y~\ref{ch:encapdecap}).

  \item \textbf{¿Qué hay dentro de una clave secreta?} En el
        conjunto mayor, $14\,120$ bytes: los $12\,800$ que codifican
        el orden de los elementos $\alpha$ (red de Benes), los
        $1\,024$ del vector de rechazo $s$, los $256$ del polinomio
        $g$ y las semillas ($40$). El orden de $\alpha$, no el
        código---que es público---, es el secreto
        (capítulo~\ref{ch:keygen}).
\end{enumerate}

\section{Ecosistema y estándares}

\begin{enumerate}[leftmargin=1.9em,itemsep=5pt,start=8]
  \item \textbf{¿Está estandarizado?} Sí, por ISO/IEC: la
        Enmienda~2 a la 18033-2 (junio de 2026) incluye Classic
        McEliece con cuatro tamaños y cuatro variantes---$16$
        perfiles---junto a ML-KEM y FrodoKEM~\cite{iso18033,cmiso}.
        El NIST \emph{no} lo estandarizó: su proceso de KEMs cerró
        en marzo de 2025 con HQC, y en 2025 aplazó la decisión sobre
        McEliece citando la estandarización ISO en
        curso~\cite{nist8545,cmnist}. Norma ISO hay; FIPS, no.

  \item \textbf{¿Dónde se usa hoy?} En despliegues ligados a la
        transición post-cuántica: túneles VPN (Mullvad), Rosenpass,
        PQConnect, y en el ecosistema de implementaciones
        liboqs/PQClean y SUPERCOP~\cite{bernstein2026,cmimpl}.
        Todos en modo híbrido.

  \item \textbf{¿Por qué el NIST lo dejó fuera?} Por ingeniería, no
        por criptoanálisis: cientos de kilobytes de clave pública no
        encajan en sus casos de uso objetivo (TLS de alto
        volumen), y en 2025 pesó la coexistencia con ISO. Es el
        patrón de toda la historia del esquema: seguridad cómoda,
        ergonomía incómoda (capítulo~\ref{ch:parametros}).

  \item \textbf{¿Qué fue «You Only Decapsulate Once»?} Una
        afirmación de 2025 sobre ataques de desencapsulación
        múltiple; el equipo respondió con su documento del 20 de
        noviembre~\cite{cm2025yodo}. Precursor del patrón de
        2025--2026: afirmación fuerte, respuesta pública,
        contabilidades en desacuerdo (capítulo~\ref{ch:2026}).
\end{enumerate}

\section{Seguridad: qué se sabe y qué no}

\begin{enumerate}[leftmargin=1.9em,itemsep=5pt,start=12]
  \item \textbf{¿Está roto?} No. A octubre de 2026 no hay ataque
        práctico contra ningún conjunto, ni recuperación de clave
        ejecutada contra una clave real, ni rotura de los retos de
        la especificación. Pero el tono cambió: aparecieron
        distinguidores que alcanzan el régimen real (GIJS, BLRT,
        Hemmert--Wiemers) y estimaciones de recuperación muy por
        debajo de ciertas cifras declaradas---de $2^{99}$--$2^{102}$
        en la contabilidad de una sola ejecución de Weis a
        $2^{140}$--$2^{148}$ en las estimaciones del BSI---frente a
        los $2^{140{,}8}$--$2^{275{,}6}$ del modelo clásico de
        memoria libre según el conjunto~\cite{ghoshal2026,weis2026,bsi2026}.

  \item \textbf{¿Qué significa «distinguir», exactamente?} Que un
        algoritmo decide si una matriz pública es de Goppa o
        aleatoria. No recupera $e$ ni la clave: es una propiedad más
        débil (figura~\ref{fig:seg-cadena}). Durante décadas los
        distinguidores vivieron en regímenes que McEliece no usa;
        los de 2026 alcanzan los parámetros reales, aunque con
        costes aún fuera de alcance y sin reducción conocida de
        distinguir a romper~\cite{ghoshal2026}.
        Capítulos~\ref{ch:estructura} y~\ref{ch:2026}.

  \item \textbf{¿Y la amenaza cuántica?} No hay rotura cuántica
        conocida: Shor no aplica (no hay grupo abeliano oculto), y
        los mejores ataques cuánticos ganan a lo sumo la mitad del
        exponente~\cite{kachigar2017}. ISO partió esos exponentes y
        exigió $\ge2^{128}$: excluyó a \texttt{mceliece348864} y
        admitió los demás con matices~\cite{iso18033,cmiso}
        (capítulo~\ref{ch:parametros}).

  \item \textbf{¿Qué dice el BSI?} Su nota del 29 de septiembre de
        2026 es la advertencia institucional más dura recibida: los
        métodos de agosto--octubre merecen tomarse en serio; «no
        debería seguir usándose en nuevos desarrollos»; los
        despliegues híbridos existentes conservan la mitad clásica;
        recomienda ML-KEM, FrodoKEM y HQC, y revisará su
        TR-02102-1~\cite{bsi2026}.

  \item \textbf{¿Y qué responde el equipo?} Que ningún ataque
        completo está verificado; que las cifras agresivas dependen
        de contabilidades discutibles (operaciones de bits, memoria
        gratis, pasos heurísticos no ejecutados); y que el suelo
        trivial de recuperación de clave es $2^{256}$, la búsqueda
        exhaustiva de la semilla de 32 bytes~\cite{cm529,cm610,bernstein2026}.
        El desacuerdo es técnico, público y está abierto
        (capítulo~\ref{ch:2026}).
\end{enumerate}

\section{Práctica: elegir, desplegar, verificar}

\begin{enumerate}[leftmargin=1.9em,itemsep=5pt,start=17]
  \item \textbf{¿Qué conjunto y qué variante elegir?} La
        recomendación del equipo es simétrica y simple: los seis
        dígitos para uso a largo plazo (\texttt{mceliece6688128},
        \texttt{mceliece6960119}); tamaños intermedios solo con
        presupuesto ajustado; \texttt{f} si el coste de generación
        importa ($<2^{-30}$ de reintento, sin efecto apreciable en
        seguridad); \texttt{pc}/\texttt{pcf} si se quiere
        confirmación de conocimiento del error ($+32$
        bytes)~\cite{cmsecurity,cmiso}. La guía completa está en el
        capítulo~\ref{ch:parametros}.

  \item \textbf{¿Cómo verifico una implementación?} Con los
        vectores de prueba de la especificación: la clave entera es
        reproducible bit a bit desde la semilla de 32 bytes, así
        que cualquiera regenera clave y criptograma y compara. Para
        C portable están PQClean y liboqs; para medidas
        comparables, SUPERCOP~\cite{cmimpl}.

  \item \textbf{¿Y si tengo que decidir hoy?} Tres reglas
        prudentes, coherentes con la nota del
        BSI~\cite{bsi2026}: (1)~en entornos con guía BSI, no
        iniciar despliegues nuevos y preferir ML-KEM, FrodoKEM o
        HQC; (2)~donde ya haya Classic McEliece, mantener el modo
        híbrido y monitorizar la literatura (la mitad clásica cubre
        la eventualidad de una degradación del componente
        post-cuántico); (3)~si el diseño exige diversidad de
        familias, tratar los tamaños \texttt{6*} como el mínimo
        razonable y revisar los resultados de 2025--2026 antes de
        congelar parámetros~\cite{cm610,cmsecurity}.
\end{enumerate}

\begin{observacion}{La respuesta honesta a «¿lo uso?»}{faq-uso}
Ningún documento de 2026---ni los ataques, ni la nota del BSI, ni
las réplicas del equipo---afirma que Classic McEliece esté roto
(octubre de 2026). Y ninguno autoriza a ignorar la tendencia. La
cifra que mejor resume el momento es la del capítulo~\ref{ch:2026}:
las estimaciones agresivas bajaron de $2^{9116}$ a $2^{94}$--$2^{102}$
en \emph{una parte} de la contabilidad, mientras la otra sigue en
$2^{142}$--$2^{150}$; la distancia entre ambas contabilidades es,
literalmente, un problema abierto. Un usuario racional hace lo que
hacen los usuarios reales: híbrido, tamaños grandes, seguimiento
activo, y cero apuestas a que el desacuerdo se resuelva a su favor.
\end{observacion}
